% Generado por generar_capitulos.py a partir de ../.tex8; no editar a mano
\chapter{Conclusiones}\label{conclusiones}
\label{sec:conclusiones}
% <- 07
Este trabajo construyó EAM-TMS, un sistema distribuido de memoria transactiva hecho de
Memorias Asociativas Entrópicas ponderadas (W-EAM) y de su extensión heteroasociativa
(EHAM)~\cite{pineda2021entropic,pineda2022weighted,morales2025missing}, y lo estudió en
cinco fases.
Un prototipo de tres módulos sirvió para comprobar que el directorio emerge de la experiencia y
para elegir cómo leerlo.
Después el sistema creció a los ocho módulos de ETH-80, con una memoria heteroasociativa
por módulo, y sobre esa versión se midieron sus propiedades y se integraron la voz y la
cámara.
Las decisiones de enrutamiento y de rechazo no usan clasificadores discriminativos; la
partición del contenido usa las etiquetas de clase, y el codificador visual se entrenó
con una pérdida de clasificación auxiliar, de la que depende el enrutamiento visual.
Del modelo de Wegner, el sistema implementa la actualización del directorio y la
coordinación de la recuperación; la asignación de información la hace la etiqueta.

Entre el prototipo y el sistema final cambiaron a la vez el codificador, la cuantización
de las pistas, las memorias por módulo, el banco de consultas y el mediador, de modo que
las cifras de tres módulos no se comparan con las de ocho como un experimento controlado.

\section{Con respecto a las preguntas de investigación}
% <- 07, «Respuesta a la pregunta de investigación» hasta antes de «Trabajo futuro»
\paragraph{Emergencia del directorio (fase 1).}
El directorio emerge de la experiencia con la arquitectura de la fase 1: tres módulos,
cuatro memorias por módulo, cuantización por signo, TME mediador y banco de 80 consultas.
Empieza vacío, y tras una sola interacción la fase madura acierta el 26.25\,\% de las
consultas y rechaza el 72.5\,\%.
A las 13 interacciones acierta el 91.25\,\%, y con las 80 del banco, el 98.75\,\%, con
un rechazo del 1.25\,\%.
La prueba más directa es el orden bloqueado por dominio.
Mientras solo se presentaron consultas de \texttt{apple}, el directorio solo enrutó a
\texttt{apple} y se quedó en el 33.75\,\%, la parte de esa clase en el banco; cruzó el
90\,\% hasta la interacción 67, ya dentro del bloque de \texttt{car}.
El orden intercalado, el barajado y el bloqueado terminan en las mismas cuentas,
[78, 65, 53].
En el directorio queda a quién vio ganar cada pista, y lo que no presenció no lo enruta.
Esta respuesta tiene cuatro límites.
La curva se midió sobre las mismas 80 consultas que forman el directorio y mide cuánto
del banco cubren los tokens ya registrados.
El techo lo pone la fase temprana (98.73\,\% acumulado frente al 98.75\,\% de la
madura).
Sin la lectura B1, la fase madura no lo alcanza.
Por último, la formación con ocho módulos no repite el umbral: ninguna condición pasa del 90\,\% y
todas terminan en el 81.2\,\% (Sección~\ref{sec:formacion}).

\paragraph{Lectura del directorio (fase 2).}
La lectura elegida en la fase 2 es B1, que divide la evidencia de cada módulo entre su
número de registros más uno.
En el prototipo, sobre las 80 consultas, la lectura cruda deja la fase madura en el
53.75\,\%, porque \texttt{apple}, el módulo con más registros, gana el 78.75\,\% de las
consultas.
B1 la sube al 98.75\,\%, y dividir entre la raíz del número de registros (B2) se queda en
el 75.0\,\%.
La elección se sostuvo con ocho módulos.
Sobre conjuntos fijos de 120 consultas que nunca forman el directorio, B1 supera a la
lectura cruda con todos los tamaños de formación y en todos los sorteos, por 4.5 a 8.8
puntos en promedio, y con ella la fase madura sube del 78.9\,\% con 50 consultas de
formación al 91.0\,\% con 291, frente al 89.7\,\% de la temprana
(Sección~\ref{sec:calibracion-fija}).
En el sistema de ocho módulos, la lectura cruda le da a \texttt{apple} el 20.5\,\% de las
consultas aunque las cuentas del directorio están casi parejas (2.974 de 3.000 bits), y B1 lo deja
en el 13.8\,\% (Sección~\ref{sec:calibracion}).
Leer el directorio con el principio de tamaño da en ETH-80, donde el registro está
balanceado, las mismas decisiones que B1 salvo una de 656 imágenes y 8 de 411 consultas.
B1 opera fuera de las memorias y corrige un sesgo de la lectura sin tocar lo registrado.

\paragraph{Escalado a ocho módulos (fase 3).}
El escalado a los ocho módulos de ETH-80 funciona con una memoria heteroasociativa y dos
directorios por módulo.
La caída del protocolo completo, del 98.75\,\% de la fase madura con tres módulos al
81.2\,\% con ocho, no se puede atribuir al número de módulos, por los cambios simultáneos
del escalado.
En la arquitectura final, el hemisferio visual enruta el 96.2\,\% de las 656 imágenes de
prueba con dos falsos enrutamientos, y 113 de las 116 imágenes enrutadas entre las 15
primeras de cada clase evocan palabras de su dominio
(Sección~\ref{sec:hemisferio-visual}).
Ante objetos que ningún módulo vio, la cobertura de punta a punta es del 89.2\,\%; con un
codificador reentrenado sin el objeto reservado es del 87.6\,\%, y los falsos
enrutamientos suben de 2 a 136 de 3280 (Sección~\ref{sec:objetos-nuevos}).
Las memorias homoasociativas que antes se llenaban aparte equivalen a lecturas de la
heteroasociativa: con una sola memoria de contenido por módulo, el sistema toma las
mismas decisiones en imagen y casi todas en texto (Sección~\ref{sec:una-memoria}).

Sobre el tamaño hay dos mediciones controladas, las dos con la arquitectura final.
Frente a una sola EHAM con las ocho clases y el mismo contenido, el sistema acierta la
clase igual desde 50 imágenes por clase; con 200, sobre las 171 consultas reservadas,
aciertan el 96.8 y el 96.9\,\% ($p = 0.86$).
La partición compra fidelidad del recuerdo (lo recordado queda 5.4 unidades más cerca de
una instancia real), menos mezcla cuando la pista es compartida (71.4 frente a 51.1\,\% de
clase correcta) y un rechazo que no empeora con el contenido (acepta 2 de 40 consultas
fuera de dominio, y la EHAM única, 8); esa ventaja en el rechazo depende de una segunda
compuerta decidida después de ver esas consultas (limitación (g)).
A cambio, en imagen responde al 81.2\,\% de las imágenes, y la EHAM única al 100\,\%
(Sección~\ref{sec:eham-unica}).
En los 28 pares de clases de ETH-80, una EHAM de dos clases recuerda a 2.54 unidades de
distancia más que el especialista, en promedio, y la de ocho a 5.66, peor en 26 de los 28
pares.
Esa brecha no sigue al solapamiento entre los dos dominios ($\rho = 0.01$), y el error de
clase sí ($\rho = 0.49$, $p = 0.008$).
Mezclar más dominios en una memoria degrada el recuerdo antes que el acierto de clase;
hasta 200 imágenes por clase, la memoria única no llega a un techo en el acierto, y dónde
lo alcanza queda por medir.

Frente a decisores simples sobre los mismos vectores, la memoria no agrega acierto de
clase.
En las 80 consultas del protocolo completo, el voto por tokens del vocabulario, el
centroide y los vecinos más cercanos aciertan entre el 76.0 y el 81.5\,\%, frente al
81.2\,\% del sistema, y en imagen un centroide o un vecino más cercano aciertan de 652 a
655 de las 656 imágenes, frente a 631 (Sección~\ref{sec:lineas-base}).
Agrega un rechazo que en texto no usa umbral, aunque equivale a la pertenencia
al vocabulario, el recuerdo del contenido asociado y decisiones que se leen en los
registros de sus memorias.
Las representaciones deciden la clase: en texto, el léxico que comparten clases vecinas
concentra los errores (\texttt{pear} $\to$ \texttt{apple} 9 y \texttt{horse} $\to$
\texttt{cow} 6 de 28), como en la sobrecarga de pistas de la memoria
humana~\cite{watkins1975buildup}, y en imagen, sin la pérdida de clase del codificador, el
enrutamiento baja al 70.4\,\% (Sección~\ref{sec:codificador-sin-clase}).

\paragraph{Propiedades del sistema escalado (fase 4).}
En el sistema escalado, el directorio funciona como un índice y, leído al revés, como un
registro de lo que sabe cada miembro.
Para enrutar solo cuenta quién ganó cada pista.
Un directorio único para texto e imagen enruta el texto igual que uno separado (96.2\,\%
en los dos), y las dos modalidades no se transfieren nada.
Leído de módulo a pista, el directorio reconstruye el dominio propio y el de los módulos
ajenos que presenció (80 de 80 y 112 de 112), y lo que devuelve es una envolvente del
dominio, de 2.9 a 5.5 veces más lejos de cualquier instancia real que las instancias
entre sí.

Registrar solo lo que cada módulo presencia le cuesta poco acierto al grupo.
En imagen, encadenar los directorios parciales iguala al registro compartido (81.9\,\%
con una de cada 64 transacciones ajenas presenciadas); en texto lo iguala sumar los
puntajes de los módulos conocidos (91.0 frente a 93.6\,\% con una de cada 32).
Cuando falta un especialista, el directorio sigue apuntando a él, de modo que el grupo
pierde la respuesta y conserva la ubicación, como en la pérdida de un miembro que
describe Wegner~\cite{wegner1987transactive}.
Sustituirlo por el segundo mejor responde el 24.5\,\% de las veces y nunca con la clase
correcta, mientras que describirlo con la lectura inversa del directorio responde siempre
y con la clase correcta.
El reconocimiento queda en el especialista, que acepta el 98.3\,\% de las
consultas de su dominio, y un módulo ajeno, el 16.1\,\%.
Estos experimentos de lectura inversa, perspectiva y pérdida de un miembro usan
protocolos anteriores al sistema final, con directorios compartidos o con un directorio
visual formado sin variantes.

Reconocer y recordar se separan dentro de una misma memoria.
La memoria acepta lo que su envolvente cubre, aunque no se parezca a nada visto; el
codificador lleva las entradas sin estructura cerca del origen, ese punto cae en el
soporte de \texttt{horse}, y el umbral de energía $\tau$ cierra 137 de 152 de esas
entradas.
Ante un daño al azar, el módulo deja de nombrar mucho antes de dejar de reconocer: destruir
el 1\,\% de las celdas deja sin nombre al 92\,\% de las imágenes, porque nombrar es una
conjunción sobre todas las coordenadas y reconocer, una suma.
Ese patrón tiene análogos en las personas, como la afasia óptica o la experiencia de la
punta de la lengua (Sección~\ref{sec:disc-una-memoria}).

De una imagen el sistema recupera el dominio con más seguridad que la palabra.
El primer término evocado pertenece al vocabulario de la clase en el 96.0\,\% de las
imágenes, y es el nombre de la clase en el 48.0\,\%.
En objetos nuevos, el color se acierta en el 57.8\,\% frente a una base del 49.5\,\%, y
un atributo visible, como el patrón o el pelo, sale con el valor mayoritario del módulo
en el 98 al 99\,\% de las vistas, de modo que el valor minoritario se acierta del 1 al
6\,\%.
La pertenencia a la clase generaliza; el atributo lo decide la densidad de lo
registrado.

El rechazo por contención falla en los dos sentidos.
De 12 consultas fuera de los ocho dominios, la memoria rechaza 10 por contención, y la
contención estricta por coordenada rechaza 23 imágenes que un vecino más cercano
clasifica bien.
El rechazo se probó con consultas de texto y con sondas visuales sintéticas; con
fotografías reales de otras categorías no se midió.

\paragraph{Voz y cámara (fase 5).}
La fase 5 es una demostración de integración, hecha con la versión v5 del sistema: ocho
módulos con tres memorias de contenido por módulo.
Una frase grabada, «Farm animal that gives milk», pasó por el detector de voz, Whisper y
el enrutamiento y llegó a \texttt{cow} en unos 0.6\,s, de los que 336\,ms fueron de
transcripción y 302\,ms de enrutamiento.
La visión analiza un cuadro en 54\,ms con la interfaz a 26 cuadros por segundo, y un
cuadro de la cámara sin ningún objeto de las ocho clases se rechaza.
En simulación, la corrección de color recupera 22, 18 y 15 de 24 imágenes de prueba sin
degradar, con degradación media y con degradación fuerte.
Son mediciones únicas de julio de 2026, sin salida cruda guardada; la frase entró desde
un archivo de audio y no se habló en vivo, y con la cámara física y hojas impresas solo
\texttt{cow} se reconoció de forma fiable.
No hay tasa de error de transcripción ni aciertos de enrutamiento con voz.
Esta fase muestra que la cadena funciona de punta a punta; cuánto acierta está por
medir.

\paragraph{Lo que se sostuvo y lo que cambió.}
La Sección~\ref{sec:creencias} contrasta cada creencia inicial con lo que se encontró.
Se sostuvieron tres decisiones: la lectura normalizada por conteo, elegida con tres
módulos y vigente con ocho; el directorio como índice para enrutar, que el directorio
unificado de texto e imagen confirmó; y la idea de que el acierto lo ponen sobre todo las
representaciones, con el matiz de que el mecanismo tiene límites propios.
La ventaja de la fase madura sobre la temprana cambió, porque venía de evaluar sobre las
consultas que formaron el directorio, y con un conjunto fijo la madura solo alcanza a la
temprana.
También cambió la coincidencia completa entre fases del primer prototipo, que contaba
dos rechazos como coincidencias.
Al pasar a ocho módulos, la cuantización por signo de las pistas se reemplazó por una por
magnitud.
De partir el contenido entre especialistas se esperaba más acierto de clase que con una
sola memoria; dio el mismo, junto con fidelidad, menos mezcla y rechazo.
Tampoco se sostuvo que todas las decisiones fueran operaciones de memoria, porque la suma
por tokens, B1, el umbral $\tau$ y el rescate léxico operan fuera de ellas
(Sección~\ref{sec:auditoria}).

\section{Con respecto al objetivo general}
El objetivo general (Sección~\ref{secc:og}) se cumplió.
El sistema se construyó y se midió contra decisores simples, contra una sola memoria con
las ocho clases y contra un codificador sin la pérdida de clase.
Lo que dejaron esas comparaciones se resume por objetivo específico.

\paragraph{Arquitectura (objetivo~\ref{oe:arquitectura}).}
Son ocho módulos, cada uno con una memoria de contenido heteroasociativa y dos
directorios, y un mediador que deja de intervenir en la fase madura.
Registrar, reconocer, recuperar y enrutar son operaciones de las memorias
(Sección~\ref{sec:operaciones}); la suma por tokens, la lectura B1, el umbral $\tau$ y el
rescate léxico operan fuera de ellas (Sección~\ref{sec:auditoria}).
Una sola memoria de contenido por módulo basta: los pesos de cada dominio se leen de la
memoria heteroasociativa, y con ellos el sistema toma las mismas decisiones que con
memorias homoasociativas separadas en imagen y casi todas en texto (cambian algunos
destinos del protocolo y una consulta de 80 en la curva de capacidad;
Sección~\ref{sec:una-memoria}).

\paragraph{Formación del directorio (objetivo~\ref{oe:formacion}).}
Sobre conjuntos de 120 consultas que nunca forman el directorio, la fase madura con la
lectura B1 sube del 78.9\,\% con 50 consultas de formación al 91.0\,\% con 291, y la
temprana es del 89.7\,\% (promedio de cinco sorteos del conjunto de evaluación).
Con unas 200 consultas, la fase madura alcanza a la temprana: la diferencia es de
$+0.2$ puntos con 200 y de $+1.3$ con 291, y ninguna diferencia positiva se sostiene
entre sorteos.
El directorio coincide, entonces, con la decisión del mediador en el 92.5\,\% de las
consultas evaluadas (sorteo principal), con un acierto equivalente y sin difundir la
consulta a todos los módulos.
La lectura B1 supera a la cruda con todos los tamaños de formación y en todos los
sorteos, por 4.5 a 8.8 puntos en promedio; sin ella, la fase madura nunca alcanza a la
temprana.

\paragraph{Frente a decisores simples y a una memoria única
(objetivo~\ref{oe:lineas-base}).}
En texto, el sistema acierta lo mismo que los decisores simples.
Sobre las 80 consultas del protocolo completo, en las que solo 11 traen el nombre de
su clase, acierta el 81.2\,\%, y el voto por tokens del vocabulario, el centroide y los
vecinos más cercanos dan entre 76.0 y 81.5\,\%.
Sobre las 171 consultas reservadas, con 200 imágenes por clase, enruta bien el
96.8\,\% y la memoria única clasifica bien el 96.9\,\% (diferencia de $-0.1$ puntos,
$p = 0.86$).
Lo que la memoria agrega sobre el voto en texto es el desempate entre las clases que
contienen algún token de la consulta: 2.9 puntos sobre las 400 consultas de
calibración (90.0 frente a 87.1\,\%).
Una tabla de conteos de tokens formada con las mismas consultas no se distingue del
directorio (92.2 frente a 93.0\,\%, $p = 0.58$).

\paragraph{Rechazo (objetivo~\ref{oe:rechazo}).}
De las 40 consultas fuera de dominio, el sistema, con la doble compuerta, acepta 2 con
25, 50, 100 y 200 imágenes por clase, y la memoria única 2, 2, 5 y 8; la segunda
compuerta se decidió después de ver estas consultas (limitación (g)).
Los decisores de coseno aceptan las 40.
La compuerta textual coincide con la pertenencia al vocabulario en los 3064 pares de
token y clase examinados, de modo que un voto léxico rechaza las mismas consultas.
Las dos que pasan son las únicas que comparten una palabra con el vocabulario de un
módulo (limitación (a)).
En imagen, el umbral de energía y la contención rechazan buena parte de las entradas sin
estructura (limitación (c)).
El rechazo de imágenes reales de otras categorías no se midió.

\paragraph{Recuperación en ambas direcciones (objetivo~\ref{oe:recuperacion}).}
De texto a imagen, el sistema responde con una imagen al 98.7\,\% de las 171 consultas
reservadas, y la memoria única al 100\,\%.
La distancia de la imagen recordada a la imagen real más cercana de su clase es de
27.2 con 200 imágenes por clase, frente a 32.5 de la memoria única y 26.8 del oráculo.
De imagen a texto, 113 de las 116 imágenes enrutadas entre las 15 primeras de cada clase
evocan palabras de su dominio (97.4\,\%).

\paragraph{Imágenes y codificador (objetivo~\ref{oe:vision}).}
Los directorios visuales envían al especialista el 96.2\,\% de las 656 imágenes de
prueba (631), rechazan 23 y se equivocan en dos.
El centroide y los vecinos más cercanos sobre los mismos códigos aciertan de 652 a 655
sin rechazar ninguna, y el vecino más cercano acierta las 23 que el sistema rechaza.
Con un codificador entrenado sin la pérdida de clase, el enrutamiento cae al 70.4\,\%,
los falsos enrutamientos suben de 2 a 169 y la cobertura de punta a punta con objetos nuevos baja del
89.2 al 59.3\,\%.
La separación entre clases la hace el codificador supervisado, y la propiedad de
fallar con un rechazo en lugar de un error depende de esa supervisión.
Con objetos que ningún módulo vio, la cobertura es del 89.2\,\% frente al 95.5\,\% con
vistas nuevas de objetos conocidos, pero el codificador sí había visto esos objetos;
con un codificador reentrenado sin el objeto reservado, la cobertura es del 87.6\,\%
[84.3, 91.0] y los falsos enrutamientos suben de 2 a 136 de 3280 (4.1\,\%).

\paragraph{En suma.}
El sistema enruta como un decisor simple y, además, rechaza lo que no contiene,
recupera en ambas direcciones y guarda un registro de quién sabe qué.
Su acierto lo ponen las representaciones: el vocabulario en texto y el codificador
supervisado en imagen.


\section{Recomendaciones para trabajo a futuro}
% <- 07, «Trabajo futuro» y párrafo final
Los siguientes pasos salen de las limitaciones de la Sección~\ref{sec:limitaciones}:
\begin{enumerate}
    \item \textbf{Partición de validación y bancos nuevos} (limitaciones (g) y (l)): elegir el
    punto de control del codificador, el número de variantes, la lectura del directorio y
    las compuertas sin usar el conjunto de prueba, y confirmar las cifras de texto, de
    imagen y fuera de dominio con bancos nuevos, escritos por terceros sin acceso al
    vocabulario, y con una partición reservada sorteada en lugar de la cola de la lista
    de cada clase.
    \item \textbf{Barrido del número de módulos con la arquitectura final}: tomar
    subconjuntos de 2 a 8 módulos de los modelos actuales, sin reentrenar ni volver a
    llenar, porque las memorias de contenido son independientes por módulo, y formar de
    nuevo los directorios.
    Separaría el efecto del número de módulos de los demás cambios del escalado; el
    codificador seguiría siendo el de ocho clases, y habría que declararlo.
    Con los tiempos de la fase temprana, el hemisferio visual y la fase madura, medidos en
    septiembre en otra computadora portátil, cada número de módulos y semilla costaría
    unos 17 minutos.
    \item \textbf{Pérdida de un miembro con el protocolo final}: repetir el experimento 9
    con directorios perspectivales, una memoria heteroasociativa por módulo y el
    directorio visual de 16 variantes, y fijar con esa corrida cuántas consultas llegan al
    módulo perdido según el número de interacciones.
    \item \textbf{Voz y cámara con cifras}: un banco de frases habladas en vivo, con su
    tasa de error de transcripción y su acierto de enrutamiento; latencias repetidas, con
    su distribución, y un conteo de aciertos con la cámara física y objetos impresos.
    Todo con la arquitectura final y guardando la salida cruda en \texttt{results/}.
    \item \textbf{Asignación de información y aprendizaje en fase madura} (limitación
    (e)): que una imagen o una palabra nueva entre por cualquier módulo y el grupo la
    registre en el especialista que indique su directorio, sin la etiqueta de clase (la
    asignación de información de Wegner); que los módulos sigan registrando después de
    formar el directorio, y que el directorio se reorganice cuando un módulo sale del
    grupo o entra en él.
    \item \textbf{Reconocer sin nombrar como consulta al directorio}
    (Sección~\ref{sec:disc-una-memoria}): cuando un módulo reconoce una imagen pero no
    puede nombrarla, que pregunte el nombre a quien lo sepa, como en el modelo de
    Wegner, en lugar de no responder.
    \item \textbf{Directorios perspectivales en texto y criterio de madurez}
    (Sección~\ref{sec:disc-wegner}): repetir la calibración y la curva de formación con
    directorios perspectivales, medir cuánto convergen entre sí los directorios de los
    módulos y que el paso a la fase madura lo decida un criterio interno, como que la
    información mutua entre dominio y ganador deje de crecer.
    \item \textbf{Fuentes de conocimiento complementarias a ConceptNet} (limitación
    (i)): medir si conocimiento alineado con el lenguaje de las consultas eleva el techo
    de las clases con léxico compartido (\texttt{pear} 76\,\%, \texttt{horse} 88\,\%).
    \item \textbf{Rechazo con imágenes reales de otras categorías}
    (Sección~\ref{sec:disc-lineas-base}): medir cuántas fotografías de objetos ajenos a
    las ocho clases aceptan los directorios visuales y las memorias de contenido, frente
    a los decisores simples sobre los mismos latentes, que no rechazan ninguna; el
    rechazo visual de este trabajo solo se probó con sondas sintéticas.
    \item \textbf{Evocación estricta} (Sección~\ref{sec:hemisferio-visual}): medir la
    evocación de dominio de las 116 imágenes con el criterio estricto, que solo cuenta
    las palabras exclusivas de la clase.
    \item \textbf{Lesión con un nombrar graduado} (limitación (j)): repetir la lesión con
    un nombrar que tolere una fracción de columnas vacías y con un daño difuso, y medir
    si la disociación se mantiene, si aparece la inversa y si las etiquetas evocadas
    derivan hacia las compartidas entre clases.
    \item \textbf{Curva del techo entrópico}: llenar la EHAM única con más clases y más
    contenido que en este trabajo, hasta encontrar dónde cae el acierto de clase, y
    comparar ahí con el reparto entre especialistas.
\end{enumerate}

\bigskip

Las Memorias Asociativas Entrópicas de Pineda y Morales bastan para construir un sistema
distribuido de memoria transactiva de ocho especialistas, bidireccional y con rechazo
explícito, sin clasificadores discriminativos en las decisiones de enrutamiento y
rechazo.
Desde el prototipo de tres módulos el directorio emerge de la experiencia, y la lectura elegida ahí
es la que usa el sistema de ocho.
El mecanismo necesita dos cosas que no salen de él, la normalización por conteo del
directorio y un código en que cada coordenada separe las clases.

\section{Difusión y productos derivados}
Este trabajo se presentó como ponencia oral en el 12.º Encuentro de Jóvenes Investigadores
del Estado de Chihuahua, con el título «Sistema distribuido de memorias heteroasociativas
entrópicas especializadas con una arquitectura de enrutamiento semántico y recuperación
asociativa multimodal basada en la teoría de memoria transactiva».

El mismo enfoque, memorias asociativas entrópicas organizadas como una memoria
transactiva, se aplicó además a la detección de fallas en piezas industriales.
Ese trabajo está en preparación para enviarse a la revista CULCYT de la Universidad
Autónoma de Ciudad Juárez.

